\begin{abstract}
Spectral adaptation changes which frequencies a continuous representation
prefers, yet its benefits vary across reconstruction tasks. We examine this
variation through a classical linear-estimation reference conditioned on an
observation operator, noise model and assumed signal covariance. Exact
scalar and matrix mismatch identities are distinguished from approximate
band diagnostics and empirical observations about trained networks.
An independent reimplementation supports mismatch-cost ranking (Spearman
0.903 over 60 aggregated cells). A multi-frame diagnostic ranks eight
practical GP priors across three noise levels (0.856), but selects the best
prior at only one level. Completed seed2/3 evaluations of LTE and SC-INR
support simple footprint ordering of readout losses in all four
architecture--seed groups. Scale marginalisation predicts ordering in four
groups and magnitude in three, but fails function similarity in all four
and does not consistently beat the simpler predictor. A prospective
36-case linear-GP test improves a fixed default by 1.105 dB
(95\% image-bootstrap interval [.726, 1.504]), but chooses the same prior
in every case, trails four-fold CV by .112 dB and costs 1.76 times as much.
The results support
a conditional account of mismatch costs and empirical readout-loss
structure, rather than a universal recovery boundary, a proven network
mechanism, or a cheap adaptive model selector.
\end{abstract}
